\documentclass[11pt,a4paper]{article}
\usepackage[T1]{fontenc}
\usepackage{lmodern,amsmath,amssymb,amsthm,booktabs,microtype}
\usepackage[margin=27mm]{geometry}
\usepackage[colorlinks=true,linkcolor=blue,urlcolor=blue]{hyperref}
\newtheorem{theorem}{Theorem}[section]
\newtheorem{lemma}[theorem]{Lemma}
\newtheorem{proposition}[theorem]{Proposition}
\newtheorem{corollary}[theorem]{Corollary}
\newcommand{\Q}{\mathbb Q}\newcommand{\Z}{\mathbb Z}
\setlength{\emergencystretch}{2em}
\title{A local obstruction to quadratic input maps}
\author{}\date{9 November 2026}
\begin{document}\maketitle
\begin{abstract}
Consider the five-contact polynomial
$r(T)=T(T-1)(T-2401)(32T+343)(2T-1)$.
No integer polynomial $q$ of degree at most two can simultaneously
reflect integrality over the rationals, take odd values at every
integer, and make $r(q(x))$ a square in $\Q_2$ for every integer $x$.
The obstruction follows from an exact classification of quadratic
integrality-reflecting maps and elementary two-adic calculations.
It concerns this particular scalar normalization and imposes no
lower bound on the degree of an arbitrary Diophantine definition.
\end{abstract}

\section{Integrality-reflecting quadratic maps}

Say that $q\in\Z[X]$ reflects rational integrality if
\[
        q(x)\in\Z\quad\Longleftrightarrow\quad x\in\Z
        \qquad(x\in\Q).
\]

\begin{proposition}\label{prop:reflection}
Let $q(X)=AX^2+BX+C$ with nonzero integer $A$.
Then $q$ reflects rational integrality if and only if
$|A|$ is squarefree and $A$ divides $B$.
\end{proposition}
\begin{proof}
Since $q(-B/A)=C$, reflection forces $A\mid B$.
Write $B=Ak$. If a prime square $p^2$ divides $A$, then
\[
 q(1/p)=A/p^2+Ak/p+C\in\Z,
\]
contrary to reflection. Thus $|A|$ is squarefree.

Conversely, suppose these two conditions hold. Write $x=u/v$
in lowest terms, with $v>0$. If $q(x)$ is integral, then
\[
                 v^2\mid A u(u+kv).
\]
Both $u$ and $u+kv$ are prime to $v$, so $v^2\mid A$.
Squarefreeness gives $v=1$.
\end{proof}

\section{Three nonsquare residue classes}

Put
\[
             r(T)=T(T-1)(T-2401)(32T+343)(2T-1).
\]
For a nonzero rational number $a$, write $v_2(a)$ for its
normalized two-adic valuation. An odd integer is a square in
$\Q_2$ precisely when it is congruent to one modulo eight;
in general the valuation must also be even.

\begin{lemma}\label{lem:residues}
For every integer $t$ in one of the following classes, $r(t)$
is a nonzero nonsquare in $\Q_2$:
\begin{center}
\begin{tabular}{ccc}
\toprule
Class of $t$ & $v_2(r(t))$ & $r(t)/2^{v_2(r(t))}\pmod8$\\
\midrule
$t\equiv5\pmod8$ & $4$ & $3$\\
$t\equiv7\pmod8$ & $2$ & $5$\\
$t\equiv9\pmod{16}$ & $6$ & $3$\\
\bottomrule
\end{tabular}
\end{center}
\end{lemma}
\begin{proof}
The two even factors differ by
\[
                  (t-1)-(t-2401)=2400=2^5\cdot75.
\]
If $t\equiv5\pmod8$, both have valuation two; their odd
parts differ by $600\equiv0\pmod8$, so their product is
one modulo eight. The other three factors have product
$5\cdot7\cdot1\equiv3\pmod8$.
If $t\equiv7\pmod8$, both even factors have valuation one;
their odd parts differ by $1200\equiv0\pmod8$.
The other factors have product $7\cdot7\cdot5\equiv5\pmod8$.

If $t\equiv9\pmod{16}$, both even factors have valuation three.
Their odd parts differ by $300\equiv4\pmod8$; the product
of those odd parts is therefore five modulo eight. The other
three factors have product $1\cdot7\cdot1\equiv7\pmod8$,
and $5\cdot7\equiv3\pmod8$.
\end{proof}

\section{The obstruction}

\begin{theorem}\label{thm:main}
There is no polynomial $q\in\Z[X]$ of degree at most two
satisfying all three conditions:
\begin{enumerate}
\item $q$ reflects rational integrality;
\item $q(x)$ is odd for every $x\in\Z$;
\item $r(q(x))$ is a square in $\Q_2$ for every $x\in\Z$.
\end{enumerate}
Here zero may be included among the squares; this convention
does not affect the result.
\end{theorem}
\begin{proof}
A constant polynomial cannot reflect integrality. A linear
polynomial $BX+C$ reflects integrality only when $B=\pm1$:
otherwise its value at $1/B$ is integral with nonintegral input.
For $B=\pm1$ its integer values have both parities. Thus assume
$q(X)=AX^2+BX+C$ with $A\ne0$.

By Proposition~\ref{prop:reflection}, $A$ is squarefree up to
sign and $B=Ak$. Odd integer values imply $C$ odd and
$A(1+k)$ even.

If $A$ is odd, then $k$ is odd. The values of $x(x+k)$ modulo
eight are all four even residues. Indeed, replacing $x$ by $x+4$
changes its value by four modulo eight, and replacing $x$ by $x+2$
changes it by two modulo four. Multiplication by the odd number
$A$ therefore makes $q(x)$ attain every odd residue modulo eight.
In particular it attains seven, contrary to
Lemma~\ref{lem:residues}.

We may therefore write $A=2a$ with $a$ odd.
If $k$ is odd, $x(x+k)$ attains both zero and two modulo four.
Thus $q(x)$ attains both $C$ and $C+4$ modulo eight. One of
these is five or seven, again a contradiction.

It remains to consider $k=2h$. Completing the square gives
\[
                 q(x)=A(x+h)^2+c,\qquad c=C-Ah^2\text{ odd}.
\]
Translation by the integer $h$ preserves the set of integer inputs,
so it is enough to consider $q(y)=Ay^2+c$.
If $c$ is five or seven modulo eight, take $y=0$.
Otherwise $c$ is one or three. Testing odd $y$ shows that
$c+A$ must also be one or three modulo eight. Since
$A\equiv2$ or $6\pmod8$, the only remaining possibilities are
\[
 (c,A)\equiv(1,2)\pmod8
 \quad\hbox{or}\quad
 (c,A)\equiv(3,6)\pmod8.
\]
In the first case, the values at $y=0$ and $y=2$ differ by
eight modulo sixteen; one is nine modulo sixteen and is excluded
by Lemma~\ref{lem:residues}.

In the second case $c+A$ is one or nine modulo sixteen.
The latter is excluded by $y=1$. In the former case,
\[
                \frac{1-c}{A}\in\Z_2^\times,
                \qquad\frac{1-c}{A}\equiv1\pmod8.
\]
There is therefore an odd $y_0\in\Z_2$ with $Ay_0^2+c=1$.
Choose an integer $y$ with $y-y_0\equiv16\pmod{32}$.
Then $v_2(y-y_0)=4$ and $v_2(y+y_0)=1$, giving
\[
              v_2(q(y)-1)=v_2(A)+4+1=6.
\]
Since $v_2(2400)=5$, it follows that
$v_2(q(y)-2401)=5$. The other three factors of $r(q(y))$
are odd, so
\[
                         v_2(r(q(y)))=11.
\]
This odd valuation is incompatible with being a two-adic square.
Every possible case has been excluded.
\end{proof}

\begin{corollary}\label{cor:sharpness}
Among integer polynomials satisfying the three conditions of
Theorem~\ref{thm:main} and $q(n)>0$ for every $n\in\Z$, the
minimum degree is four. It is attained by $q(X)=8X^4+3$.
If $q(n)>2401$ is required, $q(X)=8X^4+2403$ attains the same degree.
\end{corollary}
\begin{proof}
A nonconstant polynomial positive on all integers has even degree.
The theorem excludes degree two, and constants do not reflect
integrality. For either displayed quartic, integrality at a reduced
rational input $u/v$ implies $v^4\mid8$, and hence $v=1$.
Its values have the required positivity and are three modulo eight.
For every integer $t\equiv3\pmod8$, the two even factors of $r(t)$
have valuation one and odd parts differing by
$1200\equiv0\pmod8$. Their odd-part product is one modulo eight,
as is the product $3\cdot7\cdot5$ of the other three factors.
Thus $v_2(r(t))=2$ and $r(t)/4\equiv1\pmod8$, proving that
$r(t)$ is a square in $\Q_2$.
\end{proof}

The conclusion concerns the simultaneous requirements in the theorem.
It does not exclude quadratic input maps when parity is unrestricted,
when an additional scalar correction is used, or when a different
contact polynomial or pole construction is chosen. It therefore
identifies a restriction of the stated normalization, rather than a
minimum degree for the overall research problem.

\section{Exact finite checks}

The standard-library program
\nolinkurl{check_quadratic_input_obstruction.py}
checks all residue representatives modulo $512$ in
Lemma~\ref{lem:residues}. It also checks the four images of
$x(x+k)$ for odd $k$ modulo eight.
For the final case of the proof it examines every pair
\[
 A\equiv6\pmod8,\qquad c\equiv3\pmod8,\qquad
 A+c\equiv1\pmod{16}
\]
modulo $128$, and explicitly finds $0\le y<32$ with
\[
                     Ay^2+c-1\equiv64\pmod{128}.
\]
These $128$ cases all have the asserted valuations six, five,
and eleven. The accompanying JSON certificate lists each witness.
The checker also verifies the square residue class $t\equiv3\pmod8$
and the two quartic examples in Corollary~\ref{cor:sharpness}.
The proof is elementary and does not depend on any arithmetic
hypothesis from a defining-polynomial construction.
\end{document}
